\cvsection{Expériences professionnelles}

\begin{cventries}
  \cventry
    {Stagiaire ingénieur de recherche en IA}
    {INRIA}
    {Bordeaux, France}
    {Mars 2026 -- Septembre 2026}
    {
      \begin{cvitems}
        \item {Analyse d’un besoin des chercheurs et développement en autonomie d’une application fonctionnelle intégrant un agent IA multi-étapes pour la recherche d’articles scientifiques.}
        \item {Automatisation en Python de la collecte de métadonnées d’articles en ligne, avec exécution par tâches cron et transmission des résultats par webhooks.}
        \item {Développement de l’agent sans framework d’orchestration, avec l’API OpenAI, des prompts adaptés au cas d’usage et la gestion des appels d’outils, de la mémoire et du contexte.}
        \item {Mise en œuvre d’un pipeline RAG avec chunking, embeddings et recherche sémantique dans ChromaDB.}
        \item {Développement d’API REST Python/FastAPI documentées avec OpenAPI, intégration de MongoDB et création d’une interface React/Next.js pour les chercheurs.}
        \item {Évaluation de LLM selon la qualité des résultats et la latence, et analyse des risques liés à l’accès aux outils et fichiers lors de tests d’OpenClaw dans un environnement Docker isolé.}
        \item {Déploiement de LLM avec vLLM sur cluster HPC équipé de GPU NVIDIA, conteneurisation de composants avec Docker et mise en place d’une CI/CD sur GitHub.}
      \end{cvitems}
    }

  \cventry
    {Stagiaire ingénieur de recherche en apprentissage par renforcement}
    {Université Karunya}
    {Coimbatore, Inde}
    {Juin 2025 -- Septembre 2025}
    {
      \begin{cvitems}
        \item {Développement d’un modèle d’apprentissage par renforcement pour la navigation autonome d’un robot hospitalier, puis déploiement et tests sur Raspberry Pi.}
      \end{cvitems}
    }
\end{cventries}
